\begin{frame}{}
  \begin{block}{Definição}
    Dizemos que $p$ é uma singularidade hiperbólica quando nenhum autovalor de $A=Df(p)$ tem parte real nula.      
  \end{block}
  Lembremos a classificação em dimensão $2$.
\end{frame}

\begin{frame}{}
  Considere o sistema linear
  \begin{equation}
    x^{\prime}=Ax,\quad A\in\mathbb{R}^{n\times n}.
  \end{equation}
  Onde a solução do problema de valor inicial $x(0)=x_0$ é
  \begin{align*}
    x(t)=e^{tA}x_0,\quad\text{para todo }t\in\mathbb{R}.
  \end{align*}
  Alem disso, em geral, temos que existem $\{v_1,v_2,\cdots,v_n\}\subseteq\mathbb{R}^{n}$ vetores da base de Jordan para $A$ associados a seus autovalores.\\
  \begin{block}{Definição}
    Chamamos $E^s$ ao \textbf{espaço estável} dado por
    \begin{align*}
      E^s=span\left\{ v_i: Re\lambda_i<0 \right\}.
    \end{align*}
    Similarmente, definimos $E^u$ como o \textbf{espaço instável} dado por
    \begin{align*}
      E^{u}=span\left\{ v_i: Re\lambda_i>0 \right\}.
    \end{align*}
  \end{block}
\end{frame}

\begin{frame}{}
  \begin{block}{Proposição}
    Seja $A$ hiperbólica, então
    \begin{align*}
      \mathbb{R}^{n}=E^{s}\oplus E^u. 
    \end{align*}
  \end{block}
  Imagen
\end{frame}

\begin{frame}{}
  O quê significa isto?\\
  Se $x_0\in E^s$, então $e^{tA}x_0\to 0$ quando $t\to +\infty$.\\
  Se $x_0\in E^u$, então $e^{tA}x_0\to 0$ quando $t\to -\infty$.\\
  Alem disso, definindo $\alpha_s=\min\left\{ |Re\lambda|:Re\lambda<0 \right\}$ temos que para cada $x\in E^s$ se satisfaz que
  \begin{align*}
    \norm{e^{tA}x}\leq ce^{-\alpha_st}\norm{x},\quad t\geq 0.
  \end{align*}
  Similarmente, se $\alpha_u=\min\left\{ Re\lambda:Re\lambda>0 \right\}$, então para $y\in E^{u}$ temos que
  \begin{align*}
    \norm{e^{tA}y}\leq ce^{\alpha_ut}\norm{y},\quad t\geq 0.
  \end{align*}
\end{frame}

\begin{frame}{}
  \begin{block}{}
    Isto deixa uma nova pergunta.\\
    Se nos modificamos $Ax$ com um termo não linear pequeno,
    \begin{align*}
      x^{\prime}=Ax+r(x),\quad r(x)=O(\norm{x}),
    \end{align*}
    a estrutura de expansão e contração ainda permanece?\\
  \end{block}
\end{frame}
